\section{Construction conjointe du continu au moyen d'histoires et de prolongations compatibles}
\label{iv:continuo-conjunto}
L’arithmétique exprimée par une forme normale conserve une lecture précise de l’état qui la produit. Pour situer cette lecture dans HMT, il est nécessaire de maintenir aussi les histoires, les feuillets et les opérations qui agissent sur eux. Ce chapitre réunit ce substrat commun. Ses différentes constructions s’obtiennent à partir du même arbre d’exécutions APP–TRIT–TPK ; les découpages de l’exposition permettent d’étudier leurs propriétés sans les transformer en domaines indépendants choisis par une application ultérieure.

Le point de départ est le noyau formel exposé précédemment. Une émission admissible contient l’opération de APP, l’orientation du TRIT, le changement de feuillet, le résidu et le quotient, ainsi que la mise à jour effective de mémoire du TPK. Lorsque seul un résidu est exprimé, l’histoire d’exécution appartient toujours à l’objet de départ. Deux exécutions ayant une lecture finale identique ne sont pas identifiées pour ce seul motif.

% BEGIN REV09_SUPERVIVENCIA_PUNTO_FIJO
\subsection{Survie par élagage des prolongements}
\label{corpus9:xi:supervivencia}

La mémoire rétrospective distingue les histoires enregistrées par les
opérations APP--TRIT--TPK ; la compatibilité prospective organise leurs
prolongements admissibles. Soit \(S\) l’espace des états complets
atteignables par ces opérations. Chaque état conserve les feuillets, les résidus,
les quotients, l’orientation, la retenue, la mémoire et les données de frontière ; lorsque la
transition dépend de la profondeur, celle-ci fait partie de l’état.
Soient \(E\subseteq S\times S\) la relation effective de prolongement et
\(A\subseteq S\) l’ensemble admissible déterminé par la compatibilité de
ces registres. Ainsi, l’admissibilité se décide au moyen des opérations et
des conditions de frontière qui engendrent l’histoire.

Pour \(X\subseteq S\), nous définissons
\begin{equation}
 \operatorname{Pre}_{E}(X)
 =\{s\in S:\exists t\in X,\ (s,t)\in E\},
 \qquad
 F_{\mathrm{surv}}(X)=A\cap\operatorname{Pre}_{E}(X).
 \label{corpus9:xi:operador-poda}
\end{equation}
La suite d’élagage s’obtient par
\begin{equation}
 V_0=A,\qquad
 V_{n+1}=F_{\mathrm{surv}}(V_n),\qquad
 V_\infty=\bigcap_{n\geq0}V_n.
 \label{corpus9:xi:sucesion-poda}
\end{equation}

\begin{proposition}[Survie et plus grand point fixe]
\label{corpus9:xi:mayor-punto-fijo}
L’ensemble \(V_n\) se compose exactement des états de \(A\) qui admettent
un prolongement de \(n\) étapes dont les états appartiennent à \(A\). La
suite \((V_n)\) est décroissante. Si chaque état a un nombre fini
de successeurs admissibles, alors
\[
 V_\infty=F_{\mathrm{surv}}(V_\infty)
          =\nu F_{\mathrm{surv}},
\]
où \(\nu F_{\mathrm{surv}}\) désigne le plus grand point fixe pour l’inclusion.
Ses éléments sont précisément les états qui admettent un prolongement
infini compatible.
\end{proposition}

\begin{proof}
Pour \(n=0\), un prolongement de longueur zéro consiste en l'état lui-même de \(A\). Si la caractérisation vaut pour \(n\), un état appartient à \(V_{n+1}\) exactement lorsqu'il appartient à \(A\) et possède un successeur dans \(V_n\) ; préfixer cette arête produit un prolongement admissible de \(n+1\) étapes. Cette équivalence prouve l'affirmation par récurrence. En outre, \(V_1\subseteq V_0\), et la monotonie de \(F_{\mathrm{surv}}\) propage \(V_{n+1}\subseteq V_n\).

Soit \(s\in V_\infty\), et écrivons
\(E(s)=\{t\in S:(s,t)\in E\}\). Comme \(s\in V_{n+1}\) pour tout \(n\),
les ensembles \(E(s)\cap V_n\) sont non vides, finis et décroissants.
Leur intersection contient un élément \(t\), qui appartient à
\(V_\infty\) et vérifie \((s,t)\in E\). Par conséquent,
\(V_\infty\subseteq F_{\mathrm{surv}}(V_\infty)\).
Réciproquement, la monotonie donne
\(F_{\mathrm{surv}}(V_\infty)\subseteq V_{n+1}\) pour tout \(n\) ;
l’image est également contenue dans \(A=V_0\), de sorte qu’elle appartient à
\(V_\infty\). Cela prouve l’égalité de point fixe.

Si \(W=F_{\mathrm{surv}}(W)\), alors \(W\subseteq A=V_0\).
De \(W\subseteq V_n\), on déduit
\(W=F_{\mathrm{surv}}(W)\subseteq F_{\mathrm{surv}}(V_n)=V_{n+1}\).
La récurrence implique \(W\subseteq V_\infty\), qui prouve la maximalité.
Enfin, l’arbre des prolongements d’un état de \(V_\infty\)
est à ramification finie et possède un niveau non vide à toute profondeur.
Le lemme de König fournit une branche infinie. Dans le sens inverse, les
préfixes de tout prolongement infini situent son état initial
dans tous les \(V_n\).
\end{proof}

\paragraph{Ramification et stabilisation.}
L'hypothèse de finitude possède un contrôle explicite. Considérons une racine avec une branche terminale distincte de chaque longueur entière positive, toutes les branches étant admissibles et disjointes après la racine. La racine admet des prolongements de toute longueur finie et appartient à tous les \(V_n\) ; chacun de ses successeurs appartient à une branche de longueur bornée. La racine ne possède donc ni prolongement infini ni successeur dans \(V_\infty\). Cet arbre a une ramification infinie à la racine.
En revanche, si \(A\) est fini, chaque inclusion stricte \(V_{n+1}\subsetneq V_n\) élimine au moins un état. L'élagage atteint un point fixe après, au plus, \(|A|\) éliminations strictes de couches. Cette borne appartient au domaine fini \(A\).

\paragraph{Unicité de la continuation.}
Pour \(s\in V_\infty\), les successeurs qui conservent la survie forment
\[
 E(s)\cap V_\infty.
\]
La sélection de l’état suivant est unique exactement lorsque
\begin{equation}
 \#\bigl(E(s)\cap V_\infty\bigr)=1.
 \label{corpus9:xi:unicidad-sucesor}
\end{equation}
Un prolongement infini est unique lorsque cette condition est maintenue
à chaque état de la route. En effet, chaque successeur survivant possède
un prolongement infini par la proposition précédente ; deux successeurs
distincts produisent deux routes distinctes. Lorsque le successeur est unique à
chaque étape, les préfixes sont déterminés successivement.
La pluralité des successeurs viables exprime une multiplicité de
prolongements dans le même domaine de survie.

\begin{proposition}[Conjugaison des deux directions de prolongement]
\label{corpus9:xi:conjugacion-supervivencia}
Soit \(C:S\to S\) une involution telle que \(C(A)=A\) et
\begin{equation}
 (x,y)\in E\quad\Longleftrightarrow\quad(Cy,Cx)\in E.
 \label{corpus9:xi:involucion-prolongaciones}
\end{equation}
Notons \(V_n^+\) l’élagage pour \(E\) et \(V_n^-\) l’élagage
pour \(E^{-1}\), tous deux d’ensemble initial \(A\). Alors
\[
 C(V_n^+)=V_n^-\quad(n\geq0),\qquad
 C(V_\infty^+)=V_\infty^-.
\]
\end{proposition}

\begin{proof}
Si \(s_0,\ldots,s_n\) est une route admissible pour \(E\), la suite
\(Cs_0,\ldots,Cs_n\) est une route admissible pour \(E^{-1}\) :
pour chaque \(j\), l’hypothèse convertit
\((s_j,s_{j+1})\in E\) en \((Cs_{j+1},Cs_j)\in E\).
Tous ses états appartiennent à \(A\). La même opération appliquée une autre
fois récupère la route initiale car \(C^2=I\) ; elle établit donc une
bijection entre les prolongements finis des deux directions.
La caractérisation de \(V_n^\pm\) par les routes donne
\(C(V_n^+)=V_n^-\). Comme \(C\) est bijective, elle commute avec
l’intersection de ces ensembles, et l’on obtient l’égalité de
\(V_\infty^\pm\).
\end{proof}

L’involution transporte l’ensemble survivant d’une direction vers celui
de la direction conjuguée. La prolongeabilité infinie dans les deux
directions correspond à \(V_\infty^+\cap V_\infty^-\) ; l’unicité
d’une route conserve le critère de
\eqref{corpus9:xi:unicidad-sucesor}. La rétention d’une généalogie
et la restriction de son arbre de continuations sont deux opérations
typées sur les mêmes registres. L’attribution d’une entropie à
ces lectures exige, en outre, la mesure d’histoires correspondante.
% END REV09_SUPERVIVENCIA_PUNTO_FIJO


\subsection{Histoires admissibles et mémoire par composition}
Soit \(H_0:=\{\ast\}\). Le niveau \(H_1\) est l’ensemble fini non vide des
germes APP--TRIT admissibles, et, pour \(k\geq1\), \(H_k\) est l’ensemble des
histoires enrichies de longueur \(k\) qui admettent des prolongements TPK
compatibles à tout horizon fini. Ses éléments se construisent par émissions
successives à partir du domaine initial du noyau. Nous écrivons \(\zeta\varepsilon\)
pour le prolongement de \(\zeta\) par une émission admissible \(\varepsilon\), et
\(\rho_k(\zeta\varepsilon)=\zeta\) pour la suppression de la dernière émission. On
conserve l’étiquette produite par l’émission, même lorsque des étiquettes
différentes expriment un même état scalaire. L’espace des histoires complètes
est défini par

\[
 H_\infty^{\rm enr}:=\varprojlim_k(H_k,\rho_k).
\]
Le sous-système cofinal \(X_N:=H_{6N}\), avec \(X_0=H_0\), coïncide avec l'horloge
de sextuplets utilisée dans la preuve de cardinalité ; par cofinalité, il existe une
identification canonique
\(H_\infty^{\rm enr}\cong\varprojlim_N(X_N,\rho_{6(N+1),6N})\).

Dans les deux feuillets arithmétiques, les données locales satisfont les divisions exactes

\[
\delta_\diamond=r_\diamond+9q_\diamond,
\qquad 0\leq r_\diamond<9,
\qquad \diamond\in\{\Sigma,\Pi\}.
\]
La coordonnée orientée du TRIT admet, de manière analogue, une écriture \(\Delta=o_r+3\kappa\), avec le représentant orienté déclaré par le noyau. La conservation du quotient permet de composer ces représentations sans remplacer une émission par son résidu isolé.

Dans un système de coordonnées de mémoire, la mise à jour d’une émission s’écrit

\[
U_\varepsilon(m)=C_\varepsilon m+\kappa_\varepsilon.
\]
Pour une trajectoire \(\alpha=(\varepsilon_1,\ldots,\varepsilon_k)\), l'opération effective est la composition ordonnée de ces mises à jour. Sa partie linéaire et sa translation sont

\[
C_\alpha=C_{\varepsilon_k}\cdots C_{\varepsilon_1},
\qquad
\kappa_\alpha=
\sum_{j=1}^{k}C_{\varepsilon_k}\cdots
C_{\varepsilon_{j+1}}\kappa_{\varepsilon_j}.
\]
Le produit vide de la dernière somme est l'identité.

\begin{proposition}[Compatibilité affine de la concaténation]
\label{prop:continuo-compatibilidad-afin}
Si \(\alpha\) précède \(\beta\), alors

\[
C_{\beta\alpha}=C_\beta C_\alpha,
\qquad
\kappa_{\beta\alpha}=C_\beta\kappa_\alpha+\kappa_\beta,
\qquad
U_{\beta\alpha}=U_\beta U_\alpha.
\]
\end{proposition}
\begin{proof}
Substituer \(U_\alpha(m)\) dans \(U_\beta\) produit \(C_\beta C_\alpha m+C_\beta\kappa_\alpha+\kappa_\beta\). La décomposition en partie linéaire et translation donne les trois identités. L'associativité de la composition démontre que le résultat ne dépend pas d'une subdivision auxiliaire de la trajectoire.
\end{proof}

Cette identité de cocycle est le contenu précis de la mémoire accumulée. La quantité transportée par une seconde trajectoire dépend de la mise à jour du système de coordonnées qu’elle réalise ; ce n’est pas, en général, la somme sans transport de deux registres.

Le retour nonadique conserve une seconde distinction. Dans les coordonnées de phase et de tour, avec \(\bar r\in\{0,\ldots,8\}\), sa représentation est

\[
\sigma_9(w,r)=
\left(w+\left\lfloor\frac{\bar r+1}{9}\right\rfloor,
r+1\pmod9\right),
\qquad
\sigma_9^9(w,r)=(w+1,r).
\]
Pendant neuf étapes, il se produit exactement un franchissement de \(8\) à \(0\). La phase revient et la coordonnée de tour augmente. L'holonomie enrichie retient en outre les mises à jour de mémoire qui accompagnent ce parcours.

\subsection{Raffinement de l'arbre et mesure des histoires}
La mesure initiale est la mesure canonique du recensement : \(\mu(\ast)=1\) et \(\mu(c)=1/|H_1|\) pour chaque germe \(c\in H_1\). Dans cette section, chaque histoire possède un nombre fini et positif de prolongements. La suppression d'émissions définit les applications de raffinement. L'admissibilité du noyau fournit les prolongements ; lorsqu'une résidence utilise une émission neutre, celle-ci donne explicitement une section de \(\rho_k\). L'existence de cette section se rapporte à l'arbre qui inclut cette émission, et non à toute restriction arbitraire de routes.

Soit \(d(\zeta)\geq1\) le nombre d'émissions admissibles en \(\zeta\). L'ordre
APP et le recensement des prolongements déterminent la loi canonique
\(p_{\rm can}(\varepsilon\mid\zeta)=1/d(\zeta)\). Sans paramètres
supplémentaires, cette loi définit

\[
\mu(\zeta\varepsilon)=\frac{\mu(\zeta)}{d(\zeta)}.
\]
Une pondération non uniforme ne peut remplacer \(p_{\rm can}\) que lorsqu'elle a déjà
été produite et ajoutée au registre enrichi. Le constructeur corrélatif
de ce chapitre utilise exclusivement \(p_{\rm can}\) ; sa mesure
est donc une fonction de l'arbre APP--TRIT--TPK, et non d'un choix externe.

\begin{proposition}[Compatibilité cylindrique et espérance conditionnelle]
\label{prop:continuo-esperanza-condicional}
À chaque niveau, on a

\[
\sum_{\rho_k(y)=x}\mu(y)=\mu(x).
\]
Dans les espaces \(L^2(H_k,\mu_k)\), les applications

\[
(\mathsf J_kf)(y)=f(\rho_k(y)),
\qquad
(E_kg)(x)=\sum_{\rho_k(y)=x}\frac{\mu(y)}{\mu(x)}g(y)
\]
satisfont \(E_k=\mathsf J_k^*\), \(E_k\mathsf J_k=I\) et \(\|\mathsf J_kf\|=\|f\|\).
\end{proposition}
\begin{proof}
La première identité est la normalisation des probabilités locales. En l'appliquant à \(|f(x)|^2\), on obtient l'isométrie. En regroupant la somme par fibres de \(\rho_k\), on obtient la formule de \(E_k\) ; sur une fonction constante dans chaque fibre, cette formule restitue sa valeur.
\end{proof}

Les masses compatibles définissent une mesure de probabilité sur l'espace des histoires infinies : le cylindre d'une histoire finie reçoit sa masse \(\mu(\zeta)\). La mesure s'étend à partir de l'algèbre des cylindres, et son unicité procède du fait que ceux-ci engendrent la \(\sigma\)-algèbre. Si l'arbre est à ramification finie et que chaque histoire possède une prolongation, la limite inverse de ses niveaux par suppression des émissions est non vide. Dans la topologie cylindrique, chaque cylindre non vide a une mesure positive.

Cette mesure distingue les histoires qui convergent vers une même représentation. La mesure uniforme sur les valeurs terminales serait une lecture différente, car elle éliminerait les multiplicités que la construction vient de conserver.

\subsubsection{Lecteurs enrichis et mémoire de fibre}
\label{sec:ch-memoria-fibra}

La différence entre histoire et valeur admet un critère opérationnel. À un
niveau fini \(H_k\), soient \(q_k:H_k\to Z_k\) un lecteur et
\(M_k:H_k\to\mathcal M_k\) le registre conservé. Le lecteur enrichi
est \(\widehat q_k=(q_k,M_k)\). La mémoire distingue des états d'une même
fibre du lecteur ; sa suffisance se vérifie dans le domaine considéré.
Cette formulation intègre à l'article lui-même le développement complet de
l'information de fibre utilisée ci-après.

\begin{proposition}[Récupérabilité à un niveau fini]
\label{prop:ch-recuperabilidad-fibra}
Il existe une inverse de \(\widehat q_k\) sur son image si et seulement si \(\widehat q_k\) est injectif. Pour une variable aléatoire \(Z\) à valeurs dans \(H_k\) et de distribution \(\mu_k\), on a
\[
 \mathsf H(Z)=\mathsf H(q_k(Z))+\mathsf H(Z\mid q_k(Z)).
\]
Si \(\widehat q_k\) est injectif sur le support de \(Z\), alors
\[
 \mathsf H(Z\mid q_k(Z),M_k(Z))=0.
\]
Ici, \(\mathsf H(Z)=-\sum_z\Pr(Z=z)\log\Pr(Z=z)\), avec
\(0\log0=0\), est l'entropie finie en nats ; l'entropie conditionnelle est
la moyenne de l'entropie des distributions conditionnelles.
\end{proposition}

\begin{proof}
Un inverse récupère un unique antécédent, ce qui exige l’injectivité ;
réciproquement, celle-ci permet d’attribuer à chaque paire exprimée son unique
antécédent. Puisque \(q_k(Z)\) est une fonction déterministe de \(Z\),
regrouper la somme qui définit \(\mathsf H(Z)\) par fibres de \(q_k\) donne la
règle de la chaîne. Si la paire \((q_k(Z),M_k(Z))\) distingue le support,
chaque distribution conditionnelle correspondante est concentrée sur un seul
état et son entropie est nulle.

\end{proof}

Ne pas exprimer une coordonnée et effacer le registre qui permet de la récupérer sont
des opérations différentes. Ce résultat n’identifie pas l’entropie à la cardinalité
du continu et n’introduit pas d’interprétation thermique : son domaine, son lecteur et sa
distribution sont ceux qui viennent d’être fixés.


\subsubsection{Équivariance de la mesure sous les symétries de l'arbre}
\label{sec:ch-naturalidad-medida}

\begin{proposition}[Transport conjoint de l'arbre et de sa mesure]
\label{prop:ch-naturalidad-medida}
Soit \(g=(g_k)_{k\geq0}\) une collection de bijections \(g_k:H_k\to H_k\)
qui commute avec les troncatures, fixe la racine et transporte bijectivement
les émissions admissibles, en conservant leurs multiplicités. Alors

\[
 d(g_k\zeta)=d(\zeta),\qquad
 (g_k)_*\mu_k=\mu_k,\qquad (g_\infty)_*\mu=\mu.
\]
Plus généralement, si \(\mu_\zeta\) est la loi des continuations d'une
histoire \(\zeta\in H_k\), on a
\[
 (g_\infty)_*\mu_\zeta=\mu_{g_k\zeta}.
\]
\end{proposition}

\begin{proof}
La bijection des émissions conserve leur nombre avec multiplicités. La
distribution uniforme des germes est également invariante. Chaque trajectoire
et son image reçoivent donc le même produit de poids locaux
\(1/d(\zeta)\). On obtient l'égalité des mesures à chaque niveau et dans
chaque cylindre. La cohérence avec les troncatures définit \(g_\infty\),
et l'unicité de la mesure déterminée par les cylindres prouve l'égalité
à la limite. Le même argument, en commençant en \(\zeta\), prouve
l'affirmation sur les continuations.
\end{proof}

On obtient ainsi la naturalité de l'arbre et de la mesure : un changement cohérent
de représentation transporte les poids et n'oblige pas à les choisir à nouveau.
La conservation des émissions et des multiplicités est une hypothèse essentielle de
la proposition, et non une donnée ajoutée après la construction.


\subsection{Transport par feuillets et comptabilité exacte des sorties}
Considérons d'abord l'espace dont la base orthonormée est formée des histoires du niveau \(k\). Le transport qui conserve l'étiquette de chaque prolongement est

\[
\mathcal U_ke_\zeta=
\sum_{\varepsilon\ \mathrm{admissible\ en}\ \zeta}
\sqrt{p_{\rm can}(\varepsilon\mid\zeta)}\,e_{\zeta\varepsilon}.
\]
Les fibres successives d'antécédents distincts sont disjointes. La normalisation de \(p_{\rm can}\) démontre donc \(\mathcal U_k^*\mathcal U_k=I\). La représentation au moyen des masses des histoires est reliée à celle-ci par le changement de base diagonal qui incorpore \(\sqrt{\mu(\zeta)}\).

L'étiquette de feuillet fait partie du registre complet d'une histoire. On écrit
\[
\widehat\zeta=((\zeta_0,s_0),\ldots,(\zeta_k,s_k)),\qquad
s_j\in\{\Sigma,\Pi\},
\]
où \(s_j\) est le feuillet actif à l'étape \(j\). La prolongation ajoute l'émission et le feuillet \(s_{k+1}\) déterminé par le sélecteur opératoire du TPK, en conservant le préfixe intégralement, y compris le feuillet initial. Dans la base de ces histoires relevées,
\[
\mathcal U_ke_{\widehat\zeta}
=\sum_{\varepsilon\ \mathrm{admissible\ en}\ \widehat\zeta}
\sqrt{p_{\rm can}(\varepsilon\mid\widehat\zeta)}\,e_{\widehat\zeta\varepsilon},
\qquad
P_ke_{\widehat\zeta}=\mathbf1_{s_k=\Sigma}e_{\widehat\zeta},\quad
Q_k=I-P_k.
\]
Ainsi, \(P_k\) lit le feuillet actuel, tandis que la distinction entre histoires est maintenue par leur préfixe complet. Deux histoires initiales qui arrivent au même feuillet actif conservent des fibres successeures distinctes. La preuve d’isométrie par collections disjointes de prolongements s’applique littéralement à cette base.
L’orientation tritique est une autre coordonnée du registre : la valeur \(0\) conserve la lecture de seuil dans chaque feuillet et n’ajoute pas de troisième feuillet à cette décomposition. Nous séparons maintenant la conservation de feuillet et l’échange de feuillet :


\[
A_k=P_{k+1}\mathcal U_kP_k+Q_{k+1}\mathcal U_kQ_k,
\qquad
\eta_k=Q_{k+1}\mathcal U_kP_k+P_{k+1}\mathcal U_kQ_k.
\]
\begin{proposition}[Identité de Gram par feuillets]
\label{prop:continuo-gram-hojas}
Pour un transport linéaire \(U\), isométrique ou non, la décomposition précédente satisfait

\[
A^*A+\eta^*\eta=PU^*UP+QU^*UQ.
\]
En particulier, pour \(U=\mathcal U_k\),

\[
A_k^*A_k+\eta_k^*\eta_k=I.
\]
\end{proposition}
\begin{proof}
Dans \(A^*A\) et \(\eta^*\eta\), les termes mixtes contenant \(P_{k+1}Q_{k+1}\) s'annulent. En additionnant les termes restants apparaît \(P_kU^*(P_{k+1}+Q_{k+1})UP_k\), ainsi que son analogue avec \(Q_k\). La première identité résulte de la complémentarité des projecteurs ; la seconde utilise en outre \(U^*U=I\).
\end{proof}

La première identité est la formule générale. Remplacer son membre de droite par \(U^*U\) requiert l'annulation des blocs croisés de ce gramien. La conservation isométrique du transport des histoires fournit ici cette condition ; elle ne doit pas être attribuée automatiquement à une compression ultérieure.

Écrivons \(F_{0,0}=I\), \(F_{0,n}=A_{n-1}\cdots A_0\) et \(T_N=F_{0,N}^*F_{0,N}\).

\begin{theorem}[Décomposition finie avec résidu terminal]
\label{thm:continuo-descomposicion-terminal}
Pour chaque horizon \(N\),

\[
I=\sum_{n=0}^{N-1}
F_{0,n}^*\eta_n^*\eta_nF_{0,n}+T_N.
\]
La suite \(T_N\) est positive et décroissante, et possède une limite forte positive \(T_\infty\). La somme des sorties converge fortement vers \(I-T_\infty\).
\end{theorem}
\begin{proof}
La proposition~\ref{prop:continuo-gram-hojas} donne \(T_n-T_{n+1}=F_{0,n}^*\eta_n^*\eta_nF_{0,n}\). La somme télescopique prouve l'égalité finie et la monotonie. Pour chaque vecteur, les formes quadratiques décroissent vers une limite ; la polarisation définit une forme bornée positive et, par le théorème de représentation des formes bornées, un opérateur \(T_\infty\). La convergence est forte parce que, pour un opérateur positif \(S\leq I\), \(\|Sv\|^2\leq\langle v,Sv\rangle\) ; on l'applique à \(S=T_N-T_\infty\).
\end{proof}

Le résidu terminal est conservé jusqu'à la démonstration de son extinction dans le domaine considéré. Cette distinction empêche qu'un retour de phase, à lui seul, se transforme en une hypothèse d'épuisement de toutes les histoires.

\subsection{Bilan orienté et annulation dans le raffinement}
Une histoire porte son orientation cumulée \(o(h)\) et les décomptes de ses visites aux deux feuillets. Le bilan local correspondant est

\[
b(h)=o(h)\bigl(N_\Sigma(h)-N_\Pi(h)\bigr),
\qquad
b^+=\frac{|b|+b}{2},\quad b^-=\frac{|b|-b}{2}.
\]
Lorsqu’un lecteur exprime la moyenne sur les prolongements, on définit \(b_c=Eb_f\). L’orientation reste dans l’argument de \(b_f\) ; elle n’est pas éliminée avant le calcul de la moyenne.

\begin{proposition}[Défaut exact d'annulation]
\label{prop:continuo-defecto-cancelacion}
La quantité

\[
\kappa_{\mathrm{can}}=
\frac{E|b_f|-|Eb_f|}{2}
\]
est non négative et satisfait

\[
E(b_f^+)=(b_c)^++\kappa_{\mathrm{can}},
\qquad
E(b_f^-)=(b_c)^-+\kappa_{\mathrm{can}}.
\]
\end{proposition}
\begin{proof}
L'inégalité triangulaire pondérée donne \(|Eb_f|\leq E|b_f|\). En écrivant les parties positive et négative au moyen de \((|b|\pm b)/2\), on obtient les deux égalités.
\end{proof}

L’information omise lorsqu’on exprime seulement le bilan est l’annulation commune des contributions opposées. Cette quantité est déterminée à partir des mêmes prolongements utilisés par le transport précédent. Le facteur \(1/9\) n’apparaît que lorsque la fibre contient neuf prolongements de poids égaux ; l’espérance conditionnelle est la règle qui reste valable pour une ramification et des pondérations générales.


\subsection{Incidence ternaire des feuillets et complexe rectangulaire}
Les deux feuillets déterminent le module \(V=\mathbb F_3e_\Sigma\oplus\mathbb F_3e_\Pi\). Ses quatre directions projectives sont les droites engendrées par \(e_\Sigma,e_\Pi,e_\Sigma+e_\Pi,e_\Sigma-e_\Pi\). Elles donnent six paires non ordonnées. Les huit vecteurs non nuls sont leurs représentants orientés. Ainsi, les dénombrements \(4,6,8\) appartiennent à des opérations distinctes sur un même module de feuillets.

Sur les six positions de paires, nous définissons

\[
(Ba)_{\{L,L'\}}=a_L+a_{L'},
\qquad
s(q)=\sum_{\{L,L'\}}q_{\{L,L'\}}.
\]
Chaque direction intervient dans trois paires ; d'où \(sB=0\) dans \(\mathbb F_3\). Le lecteur \(W(q)=\mathbf1_{s(q)=0}-1/3\) est invariant sous \(q\mapsto q+Ba\). Sur l'espace ambiant uniforme \(\mathbb F_3^6\), ses moments sont

\[
\mathbb EW=0,\qquad
\mathbb EW^2=\frac29,\qquad
\mathbb EW^3=\frac2{27}.
\]
En effet, \(s\) est surjective et chaque fibre contient \(3^5\) éléments. Le lecteur vaut \(2/3\) avec la probabilité \(1/3\) et \(-1/3\) avec la probabilité \(2/3\). Ces moyennes décrivent l'espace ambiant uniforme ; la distribution induite par un sous-ensemble d'histoires doit être calculée avec ses propres multiplicités, déjà conservées dans \(H_k\).

Pour fixer les ports sans perdre les multiplicités, à la profondeur \(k\), on conserve les lectures étiquetées \(\lambda_\Sigma(\zeta)=(\zeta,\Sigma)\) et \(\lambda_\Pi(\zeta)=(\zeta,\Pi)\). Leurs ensembles sont
\[
\mathscr P_k^\Sigma=\{\lambda_\Sigma(\zeta):\zeta\in H_k\},\qquad
\mathscr P_k^\Pi=\{\lambda_\Pi(\zeta):\zeta\in H_k\}.
\]
Ils sont finis et non vides parce que \(H_k\) l’est. Les valeurs, quotients et résidus
s’expriment sur ces ports ; une égalité de valeurs n’identifie pas leurs
étiquettes. Le produit \(\mathscr P_k^\Sigma\times\mathscr P_k^\Pi\) est le
domaine de la lecture rectangulaire, au sein duquel est enregistré le support des
incidences admissibles.


Ce niveau étant fixé, nous écrivons \(I=\mathscr P_k^\Sigma\), \(J=\mathscr P_k^\Pi\) et \(A_m=\mathbb Z/3^m\mathbb Z\). L'ordre lexicographique APP détermine la branche survivante minimale et, sur celle-ci, les ports de référence \(i_0,j_0\). Aucun choix externe n'est ajouté : changer de coordonnées produit un complexe canoniquement isomorphe. Les applications

\[
\kappa(u,v)_{ij}=u_i-v_j,
\qquad
(\delta w)_{ij}=w_{ij}-w_{ij_0}-w_{i_0j}+w_{i_0j_0}
\]
définissent la suite suivante, valable sur l'anneau fini \(A_m\) :

\[
0\longrightarrow A_m\longrightarrow
A_m^I\oplus A_m^J\mathrel{\mathop{\longrightarrow}^{\kappa}}
A_m^{I\times J}\mathrel{\mathop{\longrightarrow}^{\delta}}
A_m^{(I\setminus\{i_0\})\times(J\setminus\{j_0\})}
\longrightarrow0.
\]
\begin{proposition}[Exactitude de l'incidence rectangulaire]
\label{prop:continuo-exactitud-rectangular}
La suite précédente est exacte, avec pour première flèche la diagonale constante.
\end{proposition}
\begin{proof}
L'égalité \(u_i-v_j=0\) pour chaque couple impose que toutes les coordonnées de \(u\) et de \(v\) soient une même constante. La substitution directe donne \(\delta\kappa=0\). Si \(\delta w=0\), prenons \(u_i=w_{ij_0}\) et \(v_j=w_{i_0j_0}-w_{i_0j}\). Alors \(u_i-v_j=w_{ij}\). Enfin, toute matrice du dernier terme se prolonge avec une ligne et une colonne de référence nulles ; son image par \(\delta\) est la matrice initiale. La preuve n'utilise que des sommes et des soustractions ; elle n'exige donc pas de diviser par une unité de l'anneau.
\end{proof}

La réduction de \(A_{m+1}\) à \(A_m\), en maintenant \(k,\mathscr P_k^\Sigma,\mathscr P_k^\Pi,i_0,j_0\) fixes, commute avec les deux applications. Les formules de reconstruction sont les mêmes à chaque profondeur modulaire. C'est la compatibilité arithmétique démontrée dans cette section. Le prolongement des histoires \(k\to k+1\) conserve les étiquettes de préfixe ; ses applications entre ensembles de ports sont une donnée distincte du changement d'anneau \(m\to m+1\).

Ces applications s'obtiennent par la suppression des émissions :
\[
\rho_k^\Sigma(\zeta\varepsilon,\Sigma)=(\zeta,\Sigma),\qquad
\rho_k^\Pi(\zeta\varepsilon,\Pi)=(\zeta,\Pi).
\]
La branche survivante minimale de l'ordre APP qui vient d'être construite possède des prolongements compatibles. Ses deux étiquettes définissent des ports de référence compatibles à tous les niveaux. Le transport des coordonnées est la précomposition :
\[
(L_ku)_{i'}=u_{\rho_k^\Sigma(i')},\quad
(L_kv)_{j'}=v_{\rho_k^\Pi(j')},\quad
(L_kw)_{i'j'}=w_{\rho_k^\Sigma(i'),\rho_k^\Pi(j')}.
\]
Pour le dernier module de la suite, on prolonge d'abord la matrice par zéro sur la ligne et la colonne de référence, on applique cette même formule et on restreint au nouveau complément des ports de référence. Cette règle définit \(L_k\) également lorsqu'un nouveau port a pour antécédent le port marqué.

\begin{proposition}[Naturalité du complexe rectangulaire sous prolongement]
\label{iv:rectangular-natural}
Les applications précédentes vérifient
\(L_k\kappa_k=\kappa_{k+1}L_k\) et
\(L_k\delta_k=\delta_{k+1}L_k\).
Elles commutent en outre avec la réduction des coefficients
\(A_{m+1}\to A_m\) et se composent conformément à la suppression des préfixes.
\end{proposition}
\begin{proof}
La première identité s'obtient à partir de \(u_{\rho_k^\Sigma(i')}-v_{\rho_k^\Pi(j')}\). Dans la seconde, chacun des quatre termes définissant \(\delta\) est transporté par la même précomposition. La compatibilité des ports marqués identifie leurs lignes et leurs colonnes ; lorsqu'un antécédent est le port marqué, les quatre termes s'annulent et coïncident avec le prolongement par zéro. Toutes les formules sont des sommes et des différences à coefficients entiers, de sorte qu'elles commutent avec la réduction modulaire. La composition de précompositions est la précomposition avec le préfixe composé, ce qui démontre la dernière affirmation.
\end{proof}

\subsection{Complexe de mémoire et compatibilité des incidences}
La même coordonnée accumulée de mémoire \(c_y\in\mathbb Z\), dans un lieu de résidence \(y\), définit un complexe libre entier. Sa réduction à \(A_m\) utilise le même registre entier avant de le réduire ; les profondeurs successives sont donc compatibles. Ses générateurs sont \(f_y\) en degré \(2\), \(e_y,b_y,r_y\) en degré \(1\), et \(g_y\) en degré \(0\) :

\[
\partial f_y=3c_y e_y+b_y,\qquad
\partial e_y=g_y,\qquad
\partial b_y=-3c_y g_y,\qquad
\partial r_y=0.
\]
L'identité \(\partial^2=0\) se vérifie sur \(f_y\) au moyen de \(3c_yg_y-3c_yg_y=0\), et elle est immédiate sur les autres générateurs.

Si une trajectoire \(\alpha:y\to y'\) met à jour la mémoire, son transport dans le complexe envoie les générateurs homonymes sur leurs nouveaux représentants, sauf

\[
\Psi_\alpha(b_y)=b_{y'}+3(c_{y'}-c_y)e_{y'}.
\]
\begin{proposition}[Naturalité du complexe de mémoire]
\label{prop:continuo-naturalidad-memoria}
On a \(\partial\Psi_\alpha=\Psi_\alpha\partial\) et \(\Psi_{\beta\alpha}=\Psi_\beta\Psi_\alpha\). Les applications se prolongent linéairement sur l'anneau de coefficients choisi.
\end{proposition}
\begin{proof}
Sur \(f_y\), l'image de son bord est \(3c_ye_{y'}+b_{y'}+3(c_{y'}-c_y)e_{y'} =3c_{y'}e_{y'}+b_{y'}\). Sur \(b_y\), le bord de son image est \(-3c_{y'}g_{y'}+3(c_{y'}-c_y)g_{y'}=-3c_yg_{y'}\). Les deux autres cas sont immédiats. Dans une concaténation, les différences de mémoire s'additionnent télescopiquement, ce qui prouve la seconde identité.
\end{proof}

La relation entre deux représentants d'une même résidence peut également s'écrire comme un bord explicite. Pour une coordonnée entière \(v\), réduite dans le même anneau le cas échéant, soient

\[
z^-=r_y,\qquad z^+=r_y+3c_yv e_y+v b_y,
\qquad h_*=-vf_y.
\]
Alors \(\partial h_*=z^--z^+\). Sous un transport qui met aussi à jour \(v\), la correction \(K_\alpha=(v_{y'}-v_y)f_{y'}\) satisfait \(K_{\beta\alpha}=K_\beta+\Psi_\beta K_\alpha\). Les symboles \(K_\alpha\) de cette homotopie locale ne désignent pas le sceau dodécaphasique \(K\) : ils appartiennent à un autre domaine et leur fonction est spécifiée par l'identité précédente. Leur fonction peut aussi être vérifiée dans l'égalité

\[
z^+_{y'}-\Psi_\alpha z^+_y=\partial K_\alpha.
\]
Les deux membres valent \((v_{y'}-v_y)(3c_{y'}e_{y'}+b_{y'})\).

\begin{proposition}[Homologie et mémoire du complexe local]
\label{prop:continuo-homologia-memoria}
Sur \(\mathbb Z\), ou après réduction à \(A_m\), on a

\[
H_0(\Gamma)=H_2(\Gamma)=0,
\qquad H_1(\Gamma)\cong A_m[r_y]
\]
dans la réalisation modulaire, et \(H_1(\Gamma)\cong\mathbb Z[r_y]\) dans la réalisation entière.
\end{proposition}
\Needspace{9\baselineskip}
\begin{proof}
Définissons l'homotopie de degré \(1\) par \(h(g_y)=e_y\), \(h(b_y)=f_y\), et zéro sur les autres générateurs. Soient \(p\) la projection sur le terme engendré par \(r_y\) et \(\iota\) son inclusion. La vérification sur chaque générateur donne

\[
\partial h+h\partial=I-\iota p.
\]
En particulier, sur \(b_y\), les termes \(3c_ye_y\) et \(-3c_ye_y\) s'annulent. Le complexe se rétracte par homotopie sur le module engendré par \(r_y\), concentré en degré \(1\), ce qui démontre l'affirmation.
\end{proof}

Le transport \(\Psi_\alpha\) est un cisaillement de déterminant \(1\). Ainsi, ce complexe conserve explicitement la manière dont \(c_y\) change au niveau des chaînes, mais son homologie locale ne distingue pas les valeurs de cette coordonnée. La mémoire complète demeure dans les étiquettes et dans les transports de l'état source ; elle ne se restitue pas à partir de \(H_1\) isolément. Une torsion numérique dépendant de la mémoire requiert l'assemblage et la lecture déterminantale spécifiés, outre ce complexe local.

\subsection{Assemblage avec les chemins terminaux}
Pour un horizon fini, les résidences et leurs prolongements forment une catégorie de chemins. À chaque résidence \(\nu\) est attribué le module libre \(W(\nu)=A_m[\operatorname{Path}(\nu,\mathrm{Term})]\), qui conserve chaque continuation terminale comme générateur. Une trajectoire initiale agit dans \(W\) par précomposition. Le complexe \(\Gamma(\nu)\) de la section précédente agit de manière covariante au moyen de \(\Psi\).

L'assemblage utilise les deux transports dans le même module bigradué :

\[
\mathcal B_{q,r}=
\bigoplus_{\nu_0\to\cdots\to\nu_q}
W(\nu_q)\otimes_{A_m}\Gamma_r(\nu_0).
\]
On utilise le complexe normalisé des chemins : les chaînes dégénérées contenant des identités sont omises. L'horizon fini borne alors le nombre de prolongements stricts et le degré horizontal. Le bord horizontal supprime un lieu de résidence intermédiaire en composant les flèches adjacentes. À la première extrémité, il transporte \(\Gamma\) ; à la dernière, il précompose \(W\). Avec les signes alternés, on obtient le bord des chemins \(d_{\mathrm{bar}}\). La différentielle totale est

\[
D=d_{\mathrm{bar}}+(-1)^q\partial_\Gamma.
\]
Les identités de faces annulent par paires les termes de \(d_{\mathrm{bar}}^2\). La naturalité de la proposition~\ref{prop:continuo-naturalidad-memoria} implique que \(d_{\mathrm{bar}}\) commute avec \(\partial_\Gamma\) ; le facteur \((-1)^q\) annule les termes croisés du carré. Par conséquent, \(D^2=0\). La compatibilité se démontre sur les générateurs de chemins et n'est pas introduite par une étiquette d'appartenance à un complexe.

L'évaluation terminale n'exige pas non plus de choisir une histoire particulière. Pour chaque continuation \(h\), on applique la composition effective de ses mises à jour. Le résultat conserve la paire formée par \(h\) et sa lecture, avec sa multiplicité. Si \(k<\ell<N\), les continuations se décomposent en une union disjointe de préfixes jusqu'à \(\ell\) et de leurs continuations jusqu'à \(N\). La proposition~\ref{prop:continuo-compatibilidad-afin} démontre qu'évaluer directement ou évaluer ces deux parties produit la même lecture. On obtient ainsi le carré de naturalité du terminal.

Les modules de chemins, les incidences, le transport de mémoire et leurs différentielles procèdent des mêmes prolongements. Leur assemblage ne consiste pas à ajouter à un état cinq coordonnées qui seraient conservées par définition : leurs relations sont construites par des sommes, des bords et des compositions effectives, et sont vérifiées sur leurs générateurs.

\subsection{Déterminants, limite et lecture arithmétique}
À chaque emplacement, le complexe local libre \(\Gamma\) admet sa droite déterminante graduée

\[
\operatorname{Det}\Gamma=
\bigotimes_r\left(\bigwedge^{\mathrm{rk}\Gamma_r}\Gamma_r\right)^{(-1)^r}.
\]
L'assemblage total du paragraphe précédent possède, quant à lui, les modules \(\operatorname{Tot}_d\mathcal B=\bigoplus_{q+r=d}\mathcal B_{q,r}\). À horizon fini et sur le même anneau, ils sont libres de rang fini et seul un nombre fini d'entre eux est non nul. Leur droite est
\[
\operatorname{Det}(\operatorname{Tot}\mathcal B)=
\bigotimes_{q,r}
\left(\bigwedge^{\operatorname{rk}\mathcal B_{q,r}}
\mathcal B_{q,r}\right)^{(-1)^{q+r}}.
\]
Cette formule résulte de la décomposition en somme directe à chaque degré. Elle distingue la droite du complexe local et celle de l'assemblage avec tous les chemins ; leurs bases sont ordonnées au moyen des étiquettes conservées. La réduction modulaire commute avec ces puissances extérieures parce que les modules et leurs bases proviennent des mêmes modules libres entiers à horizon fixé.
La droite est définie par les modules et leurs degrés, aussi bien sur \(\mathbb Z\) qu'après réduction à \(A_m\) ; l'exposant négatif désigne le module dual. Une torsion numérique non nulle exige en outre de spécifier le complexe sur un corps, les bases et les identifications d'homologie pertinentes. Pour calculer les facteurs de Smith, on utilise la matrice entière produite avant la réduction, et non un relèvement arbitraire d'une matrice modulaire. Un bloc carré entier de déterminant non nul est inversible sur \(\mathbb Q\) ; sa réduction est inversible sur \(A_m\) seulement lorsque ce déterminant est une unité, c'est-à-dire lorsqu'il n'est pas divisible par \(3\). Le raffinement des mémoires et des incidences permet de formuler et de vérifier la stabilité des facteurs invariants pertinents ; ces conditions ne se déduisent pas de l'existence de la droite déterminante.

Le passage à la limite conserve les systèmes projectifs compatibles de résidus, les histoires et leurs multiplicités. Les carrés prouvés dans les parties précédentes restent commutatifs parce que leurs applications sont données par les mêmes formules à chaque niveau. Pour une coordonnée archimédienne, l’étape supplémentaire consiste à exprimer des intervalles cylindriques non vides, compatibles par inclusion et de diamètre tendant vers zéro. Leurs adhérences compactes emboîtées ont une intersection réduite à un seul point : l’existence s’obtient par compacité et l’unicité par la borne du diamètre. C’est la valeur de la lecture limite. Si les intervalles sont semi-ouverts, le point peut appartenir uniquement à leurs adhérences ; la règle de représentation du bord détermine alors son écriture positionnelle. La valeur ne remplace pas l’histoire qui l’a produite. L’autre représentation conserve les incidences du même état. Les deux destinations partagent leur origine sans pour autant recevoir la même information.

Les régions coinductives de la section~\ref{sec:generacion}, le registre
dodécaphasique de la section~\ref{sec:k-registro} et l’incidence exceptionnelle
de la section~\ref{exc:seccion} s’expriment à partir de ce substrat avec leurs applications
respectives. La réalisation hilbertienne conserve ensuite les étiquettes de
phase et spécifie sa propre inclusion unitale ; la dualité et les écrans
agissent sur les codomaines qui y sont construits. Les opérations conjointes ici
réunies transmettent histoires, feuillets, orientation et mémoire à ces
représentations. L’identification d’une fonctionnelle concrète à une forme
analytique extérieure requiert, à son tour, d’écrire son application sur les
générateurs et de vérifier leurs produits scalaires : les identités de
conservation du présent chapitre ne remplacent pas cette composition.


\subsection{Portée mathématique de la construction conjointe}
Les constructions d'état, d'arbre et de mesure ; les opérations corrélatives ; et
l'assemblage terminal avec naturalité et limite sont développés intégralement
dans cet article. Dans ce chapitre, leurs identités de
bilan, d'incidence et de mémoire ont été formulées et démontrées sur les mêmes
générateurs. Les cinq constructions demeurent ainsi les opérations d'un
unique arbre. Chaque spécialisation conserve visibles ses hypothèses exactes :
isométrie du transport, loi canonique des poids, traitement du résidu
terminal et non-singularité du bloc déterminantal lorsque cela s'applique.

Les équations précédentes prouvent les cocycles de mémoire, la compensation
pondérée, l'exactitude rectangulaire et la compatibilité du complexe de
mémoire à tout horizon fini, ainsi que leur passage compatible à la limite.
Une identification analytique ultérieure doit composer son opérateur sur ces
générateurs et vérifier les propriétés supplémentaires de son codomaine ; elle ne
modifie ni ne remplace la construction conjointe démontrée ici.

\iffalse
% Formulación preliminar conservada fuera del ensamblado; sustituida por la
% versión tipada y verificada que se carga al final de este bloque.
\subsection{Émergence conjointe des cinq constructions}
\label{sec:cinco-construcciones-continuo-ch}

Les opérations précédentes sont maintenant réunies sans les transformer en cinq branches. Pour une histoire survivante \(h\in X_k\), un horizon \(N>k\) et une profondeur modulaire \(m\geq1\), soit le registre de base commun
\[
 \mathcal B_{k,N}(h):=
 \bigl(\mathscr T_{k,N}(h),
       (\varepsilon_\alpha,s_\alpha,o_\alpha,
        r_{\Sigma,\alpha},q_{\Sigma,\alpha},
        r_{\Pi,\alpha},q_{\Pi,\alpha},
        \operatorname{Carr}_\alpha,
        \kappa_\alpha,\partial_\alpha,
        \operatorname{Surv}_\alpha)_\alpha;
       \operatorname{src},\operatorname{tgt},\circ,\rho\bigr),
 \tag{C1}
 \label{eq:base-comun-registros}
\]
où \(\mathscr T_{k,N}(h)\) est le sous-arbre des prolongements de \(h\)
jusqu'à l'horizon \(N\). Chaque composante du registre est une sortie des
divisions APP, de l'orientation TRIT ou des mises à jour TPK décrites
dans les sections précédentes ; aucune ne procède d'un problème classique choisi
ensuite.

Sur cette même base, on définit les cinq objets structurés
\[
\begin{aligned}
 \mathcal F_{\rm sp}(\mathcal B_{k,N})
   &:=(\ell^2(X_j,\mu_j),\mathcal U_j,A_j,\eta_j,T_j)_{k\leq j\leq N},\\
 \mathcal F_{\rm sol}(\mathcal B_{k,N})
   &:=(b,\kappa_{\rm can},C_\alpha,\kappa_\alpha)_{\alpha\subset\mathscr T_{k,N}},\\
 \mathcal F_{\rm gau}(\mathcal B_{k,N})
   &:=(I_j,J_j,B_j,\kappa_j,\delta_j)_{k\leq j\leq N},\\
 \mathcal F_{\rm coh}^{(m)}(\mathcal B_{k,N})
   &:=(\Gamma_y\otimes A_m,\partial,\Psi_\alpha,H_\bullet)_{y,\alpha\subset\mathscr T_{k,N}},\\
 \mathcal F_{\rm det}^{(m)}(\mathcal B_{k,N})
   &:=(\operatorname{Tot}\mathcal B,\operatorname{Det}(\operatorname{Tot}\mathcal B))_{k,N,m}.
\end{aligned}
 \tag{C2}
\]
Les discriminants \({\rm sp},{\rm sol},{\rm gau},{\rm coh},{\rm det}\) nomment respectivement les lectures spectrale, solénoïdale, de jauge, cohomologique et déterminantale d'un unique objet. Ce ne sont pas des domaines d'entrée. Si \(p_T\) oublie la structure ajoutée par le lecteur \(T\) et renvoie le registre \(\mathcal B_{k,N}(h)\), nous écrivons \(\widetilde{\mathcal F}_T=(\mathcal B_{k,N},\mathcal F_T,p_T)\). La sortie corrélative est le produit fibré
\[
\begin{aligned}
 \mathcal J_{k,N}^{(m)}(h):={}&
 \widetilde{\mathcal F}_{\rm sp}
 \mathop{\times}_{\mathcal B_{k,N}(h)}
 \widetilde{\mathcal F}_{\rm sol}
 \mathop{\times}_{\mathcal B_{k,N}(h)}
 \widetilde{\mathcal F}_{\rm gau}\\
 &\mathop{\times}_{\mathcal B_{k,N}(h)}
 \widetilde{\mathcal F}_{\rm coh}^{(m)}
 \mathop{\times}_{\mathcal B_{k,N}(h)}
 \widetilde{\mathcal F}_{\rm det}^{(m)}.
\end{aligned}
 \tag{C3}
 \label{eq:constructor-correlativo}
\]
L'égalité des cinq projections sur \(\mathcal B_{k,N}(h)\) oblige à
conserver le même arbre, les mêmes arêtes et le même registre ; c'est pourquoi
\eqref{eq:constructor-correlativo} n'est pas un produit cartésien de cinq
objets assemblés rétrospectivement.

La suppression de la dernière couche de l'arbre induit
\(\operatorname{Res}_{N+1,N}\). Les propositions précédentes donnent, sur les
générateurs, les carrés
\[
\begin{gathered}
 \operatorname{Res}_{N+1,N}\mathcal F_T(\mathcal B_{k,N+1})
 =\mathcal F_T(\mathcal B_{k,N}),
 \qquad T\in\{{\rm sp},{\rm sol},{\rm gau},{\rm coh}\},\\
 \operatorname{Det}(\operatorname{Tot}\mathcal B_{k,N+1})
 \simeq
 \operatorname{Det}(\operatorname{Tot}\mathcal K_{N+1})
 \otimes
 \operatorname{Det}(\operatorname{Tot}\mathcal B_{k,N}),\\
 (\partial\Psi_\alpha=\Psi_\alpha\partial),\qquad
 (L_k\kappa_k=\kappa_{k+1}L_k),\qquad
 (L_k\delta_k=\delta_{k+1}L_k).
\end{gathered}
 \tag{C4}
 \label{eq:naturalidad-D}
\]
Le changement \(A_{m+1}\to A_m\) commute également avec \(\partial\), \(\Psi\),
\(\kappa\), \(\delta\) et les droites déterminantes. On obtient donc
un unique prosystème dans \((N,m)\).

Soit \(\mathcal X_\chi\) la limite des histoires enrichies compatibles. Pour \(\xi\in\mathcal X_\chi\), désignons par \(\mathfrak L_N(\xi)\) le préfixe ordonné de ses \(N\) émissions complètes. Les fibres conjointes d'horizon sont
\[
 \mathfrak J_{N,\bullet}^{(m)}
 :=\coprod_{h\in X_0}
 \{(\mathfrak L,J):\mathfrak L\text{ est une branche de longueur }N,
       \ J=\mathcal J_{0,N}^{(m)}(h)\}.
 \tag{C5}
 \label{eq:fibras-conjuntas-horizonte}
\]
La marque \(\mathfrak L\) ne retient pas une copie tautologique de \(\xi\) : elle est
formée des émissions et des mises à jour produites sur chaque arête. Le
générateur conjoint est
\[
\begin{aligned}
 \operatorname{Gen}_{\rm cont}:\mathcal X_\chi
   &\longrightarrow\varprojlim_{N,m}\mathfrak J_{N,\bullet}^{(m)},\\
 \xi&\longmapsto
   (\mathfrak L_N(\xi),\mathcal J_{0,N}^{(m)}(h_0))_{N,m},\\
 \mathcal C_{\rm cont}^{\rm disc}
   &:=\operatorname{im}(\operatorname{Gen}_{\rm cont}).
\end{aligned}
 \tag{C6}
 \label{eq:continuo-discreto}
\]

\begin{theorem}[Émergence corrélative des cinq constructions]
\label{thm:correlacion-cinco}
Pour toute histoire survivante \(h\in X_k\), tout horizon
\(N\geq k+1\) et toute profondeur \(m\geq1\), le produit fibré
\eqref{eq:constructor-correlativo} existe et est non vide. Ses cinq
composantes émergent du même registre APP--TRIT--TPK, satisfont les
carrés de naturalité \eqref{eq:naturalidad-D}, les changements de base et le
passage à la limite. L'image \(\mathcal C_{\rm cont}^{\rm disc}\) contient cinq
aspects inséparables du même continu, et non cinq continus ni cinq
applications indépendantes.
\end{theorem}

\begin{proof}
La prolongeabilité de l'arbre fournit au moins une arête à chaque sommet intérieur. Le registre de cette arête détermine le transport isométrique et sa décomposition par feuillets, le cocycle orienté, les ports d'incidence, le complexe de mémoire et sa droite déterminante. On construit ainsi successivement \(\mathcal F_{\rm sp}\), \(\mathcal F_{\rm sol}\), \(\mathcal F_{\rm gau}\), \(\mathcal F_{\rm coh}^{(m)}\) et \(\mathcal F_{\rm det}^{(m)}\) sur la même base, ce qui prouve la non-vacuité de \eqref{eq:constructor-correlativo}. Aucune des cinq constructions n'apparaît comme coordonnée d'entrée.

La proposition de compatibilité affine prouve la composition de la mémoire ;
l'identité de Gram, l'exactitude rectangulaire et la naturalité du complexe
prouvent les quatre premières lignes de \eqref{eq:naturalidad-D}. La
décomposition de la bicomplexité selon la dernière couche prouve l'identité
multiplicative de la droite déterminant. Toutes ces opérations se calculent
sur les mêmes générateurs et commutent avec la réduction modulaire. La fibre,
la relation, le nombre, l'orientation, la retenue, la mémoire, la frontière et
la route demeurent dans \(\mathfrak L_N(\xi)\) ; la multisection, les terminaux
et les lecteurs linéaire et probabiliste demeurent dans
\(\mathcal J_{0,N}^{(m)}(h_0)\). Par compatibilité, la limite produit
\eqref{eq:continuo-discreto} sans disjonction.
\end{proof}

La reconnaissance ultérieure de l'unique continu construit au moyen de ses cinq
réalisations spectrale, solénoïdale, de jauge, cohomologique et déterminantale
appartient au langage conventionnel postérieur. Elle ne sélectionne pas ses états et
n'intervient pas dans sa génération.
\fi

\subsection{Émergence conjointe des cinq constructions}
\label{sec:cinco-construcciones-continuo-ch}

Pour \(0\leq k\leq N\), soit
\[
 \rho_{N,k}:=\rho_k\circ\rho_{k+1}\circ\cdots\circ\rho_{N-1},
 \qquad \rho_{k,k}:=\operatorname{id}_{H_k},
\]
et, pour une histoire survivante \(\zeta\in H_k\), définissons
\[
 \operatorname{Desc}_{k,N}(\zeta)
 :=\{x\in H_N:\rho_{N,k}(x)=\zeta\}.
\]
L'arbre fini \(\mathscr T_{k,N}(\zeta)\) a pour sommets
\[
 \coprod_{j=k}^{N}\{y\in H_j:\rho_{j,k}(y)=\zeta\}
\]
et pour arêtes uniquement les prolongements TPK \(a:y\to z\) avec \(z\in\rho_j^{-1}(y)\). Ainsi, ni ses sommets ni ses prolongements ne proviennent d'un arbre auxiliaire.

Pour \(x\in\operatorname{Desc}_{k,N}(\zeta)\), la branche marquée \(\lambda_{\zeta,x}\) est la suite
\[
 \zeta=x_k\longrightarrow x_{k+1}\longrightarrow\cdots
 \longrightarrow x_N=x
\]
déterminée par les troncatures de \(x\). Le registre de base commun est
\[
\begin{aligned}
 \mathscr R_{k,N}(\zeta;x):=
 \bigl(&\mathscr T_{k,N}(\zeta),\lambda_{\zeta,x};\\
 &\bigl(\varepsilon_a,s_a,o_a,
   (\delta_{\diamond,a},r_{\diamond,a},q_{\diamond,a})
       _{\diamond\in\{\Sigma,\Pi\}},
   \operatorname{Carr}_a,C_a,t_a,\operatorname{Fr}_a,
   \operatorname{Ruta}_a,\operatorname{Mem}_a,
   \operatorname{Surv}_a\bigr)_a;\\
 &\operatorname{src},\operatorname{tgt},\circ,
   (\rho_j)_{k\leq j<N}\bigr),
\end{aligned}
\tag{C1}
\label{eq:base-comun-registros}
\]
où
\[
 U_a(m_0)=C_am_0+t_a,
 \qquad
 \delta_{\diamond,a}=r_{\diamond,a}+9q_{\diamond,a}.
\]
Le symbole \(m_0\) de cette formule est une coordonnée de mémoire ; ce n'est pas la profondeur modulaire \(m\) utilisée plus bas. Le véritable espace de base est
\[
 \mathfrak R_{k,N}(\zeta)
 :=\{\mathscr R_{k,N}(\zeta;x):
       x\in\operatorname{Desc}_{k,N}(\zeta)\}.
\]
La branche marquée n'est pas une entrée ajoutée : c'est la suite des arêtes
effectivement produites qui se termine en \(x\). Le reste du registre conserve
simultanément les fibres latérales de l'arbre ; la marque ne remplace donc pas
la ramification par une trajectoire isolée.

\Needspace{15\baselineskip}
Posons \(R_m:=\mathbb Z/3^m\mathbb Z\), et soit
\(V_j(\zeta):=\operatorname{Desc}_{k,j}(\zeta)\). Le transport, les
projecteurs de feuillet et le résidu terminal se restreignent au sous-arbre par
\[
 \mathcal U_j^\zeta:=
 \mathcal U_j\!\upharpoonright_{\ell^2(V_j(\zeta),\mu_j)},
 \qquad P_j^\zeta:=P_j\!\upharpoonright_{V_j(\zeta)},
 \qquad Q_j^\zeta:=I-P_j^\zeta,
\]
\[
 A_j^\zeta:=P_{j+1}^\zeta\mathcal U_j^\zeta P_j^\zeta
             +Q_{j+1}^\zeta\mathcal U_j^\zeta Q_j^\zeta,
 \qquad
 \eta_j^\zeta:=Q_{j+1}^\zeta\mathcal U_j^\zeta P_j^\zeta
             +P_{j+1}^\zeta\mathcal U_j^\zeta Q_j^\zeta.
\]
Pour \(k\leq n\leq N\), soient
\[
 F_{k,k}^\zeta:=I,
 \qquad
 F_{k,n}^\zeta:=A_{n-1}^\zeta\cdots A_k^\zeta,
 \qquad
 T_{k,n}^\zeta:=(F_{k,n}^\zeta)^*F_{k,n}^\zeta.
\]
Les ports propres du même sous-arbre sont
\[
 \mathscr P_j^\Sigma(\zeta):=
 \{(y,\Sigma):y\in V_j(\zeta)\},
 \qquad
 \mathscr P_j^\Pi(\zeta):=
 \{(y,\Pi):y\in V_j(\zeta)\};
\]
nous désignons par \(B_j^\zeta,\kappa_j^\zeta,\delta_j^\zeta\) les restrictions de l'incidence et du complexe rectangulaire à ces ports, et par \(L_j^\zeta\) la précomposition induite par la suppression de préfixes.

Pour définir sans abréviation le lecteur de chemins, soit
\[
 W_\ell^\zeta(y):=
 R_m[\operatorname{Path}_{\mathscr T_{k,\ell}(\zeta)}
          (y,\operatorname{Term}_\ell)]
\]
et posons
\[
 \mathscr D_{q,r;k,\ell}^{(m)}[\mathscr R]
 :=\bigoplus_{y_0\to\cdots\to y_q\ \text{en}\
                  \mathscr T_{k,\ell}(\zeta)}
 W_\ell^\zeta(y_q)\otimes_{R_m}
 \bigl(\Gamma_r(y_0)\otimes_{\mathbb Z}R_m\bigr),
 \qquad
 \operatorname{Tot}_d\mathscr D
 :=\bigoplus_{q+r=d}\mathscr D_{q,r;k,\ell}^{(m)}[\mathscr R].
\]
Les formules déjà démontrées produisent, sur chaque
\(\mathscr R\in\mathfrak R_{k,N}(\zeta)\), les cinq objets
\[
\begin{aligned}
 \mathcal F_{\rm sp}[\mathscr R]
 &:=(\ell^2(V_j(\zeta),\mu_j),
       \mathcal U_j^\zeta,A_j^\zeta,\eta_j^\zeta,
       T_{k,j}^\zeta)_{k\leq j\leq N},\\
 \mathcal F_{\rm sol}[\mathscr R]
 &:=(b,\kappa_{\rm can},C_\alpha,\kappa_\alpha)
      _{\alpha\subset\mathscr T_{k,N}(\zeta)},\\
 \mathcal F_{\rm gau}^{(m)}[\mathscr R]
 &:=(\mathscr P_j^\Sigma(\zeta),\mathscr P_j^\Pi(\zeta),
       B_j^\zeta,\kappa_j^\zeta,\delta_j^\zeta;R_m)_{k\leq j\leq N},\\
 \mathcal F_{\rm coh}^{(m)}[\mathscr R]
 &:=(\Gamma_y\otimes R_m,\partial,\Psi_\alpha,H_\bullet)
      _{y,\alpha\subset\mathscr T_{k,N}(\zeta)},\\
 \mathcal F_{\rm det}^{(m)}[\mathscr R]
 &:=(\operatorname{Det}(\operatorname{Tot}
       \mathscr D_{\bullet,\bullet;k,\ell}^{(m)}[\mathscr R]))
       _{k\leq\ell\leq N}.
\end{aligned}
\tag{C2}
\label{eq:cinco-lectores-comunes}
\]
Ici, \(\mathscr D_{\bullet,\bullet;k,\ell}^{(m)}[\mathscr R]\) est le
bicomplexe normalisé
de chemins et de mémoire des sections précédentes, restreint au sous-arbre
d’horizon \(\ell\) et réduit sur \(R_m\). Le dernier lecteur conserve la
collection ordonnée de droites déterminantales jusqu’à l’horizon \(N\) ; de cette
manière, sa troncature est une projection canonique et ne présuppose pas de
factorisation déterminantale inexistante. La loi \(\mu\) est la loi canonique
du recensement, et les ports sont fixés au moyen de la branche lexicographiquement minimale
de l’ordre APP. Par conséquent, aucun des cinq lecteurs ne reçoit de pondération
ou de port extérieurs.


Pour
\(T\in\mathfrak T:=\{{\rm sp},{\rm sol},{\rm gau},{\rm coh},{\rm det}\}\),
définissons l'espace total typé
\[
 \widetilde{\mathcal F}_{T;k,N}^{(m)}(\zeta)
 :=\coprod_{\mathscr R\in\mathfrak R_{k,N}(\zeta)}
   \{\mathcal F_T^{(m)}[\mathscr R]\}_{\mathscr R},
 \qquad
 p_T:\widetilde{\mathcal F}_{T;k,N}^{(m)}(\zeta)
      \longrightarrow\mathfrak R_{k,N}(\zeta),
\]
où \(\mathcal F_{\rm sp}^{(m)}:=\mathcal F_{\rm sp}\) et \(\mathcal F_{\rm sol}^{(m)}:=\mathcal F_{\rm sol}\) ; autrement dit, l'exposant \(m\) est inerte pour ces deux lecteurs. La projection \(p_T\) oublie uniquement la structure du lecteur et conserve son registre source. Le produit fibré correct est
\[
\begin{aligned}
 \mathcal J_{k,N}^{(m)}(\zeta):={}&
 \widetilde{\mathcal F}_{\rm sp;k,N}^{(m)}(\zeta)
 \mathop{\times}_{\mathfrak R_{k,N}(\zeta)}
 \widetilde{\mathcal F}_{\rm sol;k,N}^{(m)}(\zeta)
 \mathop{\times}_{\mathfrak R_{k,N}(\zeta)}
 \widetilde{\mathcal F}_{\rm gau;k,N}^{(m)}(\zeta)\\
 &\mathop{\times}_{\mathfrak R_{k,N}(\zeta)}
 \widetilde{\mathcal F}_{\rm coh;k,N}^{(m)}(\zeta)
 \mathop{\times}_{\mathfrak R_{k,N}(\zeta)}
 \widetilde{\mathcal F}_{\rm det;k,N}^{(m)}(\zeta).
\end{aligned}
\tag{C3}
\label{eq:constructor-correlativo}
\]
L'égalité des cinq projections oblige à conserver le même arbre, les
mêmes arêtes, la même branche produite et le même registre. Pour chaque
\(\mathscr R\), les cinq constructions déterminent l'élément
\[
 s_{k,N}^{(m)}(\mathscr R)
 :=(\mathcal F_T^{(m)}[\mathscr R])_{T\in\mathfrak T}
 \in\mathcal J_{k,N}^{(m)}(\zeta).
\]
L'existence de cet élément provient des opérations de \eqref{eq:cinco-lectores-comunes}, et non de la notation du produit fibré.

La suppression de la dernière couche induit
\[
 \operatorname{res}^{\mathfrak R}_{N+1,N}
 \bigl(\mathscr R_{k,N+1}(\zeta;x)\bigr)
 :=\mathscr R_{k,N}(\zeta;\rho_N(x)).
\]
La restriction de chaque opérateur aux générateurs retenus définit
\[
 r^m_{N+1,N}:\mathcal J_{k,N+1}^{(m)}(\zeta)
 \longrightarrow\mathcal J_{k,N}^{(m)}(\zeta).
\]
Dans le lecteur déterminantal, \(r^m_{N+1,N}\) supprime uniquement la dernière
coordonnée de la collection de droites de
\eqref{eq:cinco-lectores-comunes}. La réduction
\(R_{m+1}\twoheadrightarrow R_m\) définit, par changement de base,

\[
 q^N_{m+1,m}:\mathcal J_{k,N}^{(m+1)}(\zeta)
 \longrightarrow\mathcal J_{k,N}^{(m)}(\zeta).
\]
Sur chaque générateur, on a
\[
\begin{gathered}
 r^m_{N+1,N}q^{N+1}_{m+1,m}
 =q^N_{m+1,m}r^{m+1}_{N+1,N},\\
 \partial\Psi_\alpha=\Psi_\alpha\partial,
 \qquad L_j^\zeta\kappa_j^\zeta
       =\kappa_{j+1}^\zeta L_j^\zeta,
 \qquad L_j^\zeta\delta_j^\zeta
       =\delta_{j+1}^\zeta L_j^\zeta,
\end{gathered}
\tag{C4}
\label{eq:naturalidad-D}
\]
et les compositions des applications \(r\) et \(q\) satisfont les identités
fonctorielles. La première égalité se vérifie parce que la suppression de la dernière couche et
la réduction des coefficients agissent sur des coordonnées distinctes. Les trois autres
sont les identités de naturalité prouvées précédemment. On obtient un unique
prosystème dans \((N,m)\), et non cinq limites juxtaposées.

À la racine, on définit
\[
 \mathfrak J_{N,\bullet}^{(m)}:=\mathcal J_{0,N}^{(m)}(\ast).
\tag{C5}
\label{eq:fibras-conjuntas-horizonte}
\]
La structure discrète conjointe du continu est la limite de ces objets corrélatifs :
\[
 \boxed{
 \mathcal C_{\rm cont}^{\rm disc}
 :=\varprojlim_{\substack{N\geq1\\m\geq1}}
   \bigl(\mathfrak J_{N,\bullet}^{(m)},
         r^m_{N+1,N},q^N_{m+1,m}\bigr).}
\tag{C6}
\label{eq:continuo-discreto}
\]
C'est une définition par le prosystème engendré par l'arbre complet ; ce n'est
ni l'image ni le graphe d'une application qui retiendrait l'état d'entrée comme
première coordonnée.

Les registres possèdent une marque terminale produite. Leurs projections compatibles induisent
\[
 \Pi:\mathcal C_{\rm cont}^{\rm disc}\longrightarrow H_\infty^{\rm enr}.
\]
Réciproquement, une histoire complète détermine, par ses préfixes déjà émis, les sections compatibles \(s_{0,N}^{(m)}\). On obtient ainsi, après avoir défini le continu,
\[
 \operatorname{Gen}_{\rm cont}:H_\infty^{\rm enr}
 \longrightarrow\mathcal C_{\rm cont}^{\rm disc},
 \qquad
 \Pi\circ\operatorname{Gen}_{\rm cont}
 =\operatorname{id}_{H_\infty^{\rm enr}}.
\tag{C7}
\label{eq:seccion-historias-continuo}
\]
Cette section démontre la récupérabilité ; elle n'est pas utilisée pour définir
\(\mathcal C_{\rm cont}^{\rm disc}\).

\begin{theorem}[Émergence corrélative des cinq constructions]
\label{thm:correlacion-cinco}
Pour toute histoire survivante \(\zeta\in H_k\), tout \(N>k\), tout
\(m\geq1\) et tout \(\mathscr R\in\mathfrak R_{k,N}(\zeta)\), les cinq
constructions de \eqref{eq:cinco-lectores-comunes} sont définies sur le
même registre, et \eqref{eq:constructor-correlativo} contient
\(s_{k,N}^{(m)}(\mathscr R)\). Leurs troncatures, réductions modulaires et
transports satisfont \eqref{eq:naturalidad-D}. De plus,
\(\mathcal C_{\rm cont}^{\rm disc}\neq\varnothing\), l'application \(\Pi\)
est surjective et scindée par \(\operatorname{Gen}_{\rm cont}\). Les
cinq composantes sont des aspects inséparables du même continu, et non cinq
domaines ni cinq applications indépendantes.
\end{theorem}

\begin{proof}
La survie qui définit \(H_k\) et la loi de prolongation TPK donnent
\(\rho_k^{-1}(\zeta)\neq\varnothing\) pour chaque sommet. Comme les fibres sont
finies, le lemme de Kőnig produit \(H_\infty^{\rm enr}\neq\varnothing\).

Un registre \(\mathscr R\) étant fixé, le transport des histoires et des feuillets produit \(\mathcal F_{\rm sp}[\mathscr R]\) ; la composition affine et le défaut exact d'annulation produisent \(\mathcal F_{\rm sol}[\mathscr R]\) ; les incidences et la suite rectangulaire produisent \(\mathcal F_{\rm gau}^{(m)}[\mathscr R]\) ; le complexe de mémoire et son transport produisent \(\mathcal F_{\rm coh}^{(m)}[\mathscr R]\) ; et les droites de l'assemblage normalisé de chemins produisent \(\mathcal F_{\rm det}^{(m)}[\mathscr R]\). Chaque construction agit sur les mêmes arêtes générées, et non sur un sous-domaine supposé ni sur l'espace ambiant abstrait de \(729\) mots.

Les compatibilités affine, de transport, d’incidence et de mémoire prouvent
\eqref{eq:naturalidad-D}. Tronquer la collection déterminantale de
\eqref{eq:cinco-lectores-comunes} commute par
définition avec \(r\) ; son changement de base commute avec \(q\). Les sections
\(s_{0,N}^{(m)}\) sont compatibles, de sorte qu’elles produisent
\(\operatorname{Gen}_{\rm cont}\). L’identité
\eqref{eq:seccion-historias-continuo} prouve à la fois la non-
vacuité de la limite, la surjectivité de \(\Pi\) et la récupérabilité de chaque
histoire sans l’introduire comme coordonnée de définition.


La fibre, la relation, le nombre, l'orientation, la retenue, la mémoire, la
frontière et la route demeurent dans le registre commun ; la multisection, les
terminaux et les lecteurs linéaire et probabiliste demeurent dans les cinq
constructions corrélatives. Le passage à la limite conserve donc la clôture des
dépendances et ne disjoint pas le continu.
\end{proof}

La reconnaissance ultérieure de l'unique continu construit au moyen du langage spectral, solénoïdal, de jauge, cohomologique ou déterminantal appartient à la comparaison conventionnelle ultérieure. Elle ne sélectionne pas ses états et n'intervient pas dans sa génération APP--TRIT--TPK.

